\documentclass[11pt,a4paper]{article}
% lesson-preamble.tex -- Shared preamble for all Phase 00 lessons
% Usage: % lesson-preamble.tex -- Shared preamble for all Phase 00 lessons
% Usage: % lesson-preamble.tex -- Shared preamble for all Phase 00 lessons
% Usage: \input{../lesson-preamble} (from subdomain directories)

% ---- Encoding and fonts ----
\usepackage[utf8]{inputenc}
\usepackage[T1]{fontenc}

% ---- Mathematics ----
\usepackage{amsmath,amssymb,amsthm}
\usepackage{mathtools}

% ---- Page geometry ----
\usepackage[margin=1in,headheight=14pt]{geometry}

% ---- Hyperlinks ----
\usepackage[colorlinks=true,linkcolor=red,citecolor=blue,urlcolor=blue]{hyperref}

% ---- Lists ----
\usepackage{enumitem}

% ---- Colors and boxes ----
\usepackage{xcolor}
\usepackage[theorems,skins,breakable]{tcolorbox}

% ---- Header/Footer ----
\usepackage{fancyhdr}
\renewcommand{\headrulewidth}{0.4pt}

% ---- Tables ----
\usepackage{booktabs}

% ---- Bibliography ----
\usepackage{natbib}
\bibliographystyle{plainnat}

% --------------------------------------------------------------------------
% Theorem environments
% --------------------------------------------------------------------------
\theoremstyle{plain}
\newtheorem{theorem}{Theorem}[section]
\newtheorem{lemma}[theorem]{Lemma}
\newtheorem{proposition}[theorem]{Proposition}
\newtheorem{corollary}[theorem]{Corollary}

\theoremstyle{definition}
\newtheorem{definition}[theorem]{Definition}
\newtheorem{example}[theorem]{Example}
\newtheorem{exercise}[theorem]{Exercise}
\newtheorem{assumption}[theorem]{Assumption}

\theoremstyle{remark}
\newtheorem{remark}[theorem]{Remark}

% --------------------------------------------------------------------------
% Colored boxes
% --------------------------------------------------------------------------
\newtcolorbox{keyresult}[1][]{%
  colback=green!5!white,
  colframe=green!50!black,
  fonttitle=\bfseries,
  title={Key Result},
  breakable,
  #1
}

\newtcolorbox{rlconnection}[1][]{%
  colback=blue!5!white,
  colframe=blue!50!black,
  fonttitle=\bfseries,
  title={RL/ML Connection},
  breakable,
  #1
}

\newtcolorbox{warning}[1][]{%
  colback=red!5!white,
  colframe=red!50!black,
  fonttitle=\bfseries,
  title={Warning},
  breakable,
  #1
}

% --------------------------------------------------------------------------
% Custom commands -- number sets
% --------------------------------------------------------------------------
\newcommand{\R}{\mathbb{R}}
\newcommand{\N}{\mathbb{N}}
\newcommand{\Q}{\mathbb{Q}}
\newcommand{\Z}{\mathbb{Z}}
\newcommand{\C}{\mathbb{C}}
\newcommand{\F}{\mathbb{F}}

% --------------------------------------------------------------------------
% Custom commands -- probability and expectation
% --------------------------------------------------------------------------
\newcommand{\E}{\mathbb{E}}
\newcommand{\Prob}{\mathbb{P}}
\newcommand{\Var}{\operatorname{Var}}
\newcommand{\Cov}{\operatorname{Cov}}
\newcommand{\indicator}{\mathbf{1}}

% --------------------------------------------------------------------------
% Custom commands -- convergence arrows
% --------------------------------------------------------------------------
\newcommand{\cas}{\xrightarrow{\text{a.s.}}}
\newcommand{\cprob}{\xrightarrow{\Prob}}
\newcommand{\clp}[1]{\xrightarrow{L^{#1}}}
\newcommand{\cdist}{\xrightarrow{d}}

% --------------------------------------------------------------------------
% Custom commands -- common operators
% --------------------------------------------------------------------------
\DeclareMathOperator{\dom}{dom}
\DeclareMathOperator{\epi}{epi}
\DeclareMathOperator{\conv}{conv}
\DeclareMathOperator{\relint}{relint}
\DeclareMathOperator{\aff}{aff}
\DeclareMathOperator{\cl}{cl}
\DeclareMathOperator{\interior}{int}
\DeclareMathOperator{\tr}{tr}
\DeclareMathOperator{\diag}{diag}
\DeclareMathOperator{\rank}{rank}
\DeclareMathOperator{\sgn}{sgn}
\DeclareMathOperator{\supp}{supp}
\DeclareMathOperator{\argmin}{arg\,min}
\DeclareMathOperator{\argmax}{arg\,max}
\DeclareMathOperator{\prox}{prox}
\DeclareMathOperator{\dist}{dist}
\DeclareMathOperator{\diam}{diam}
\DeclareMathOperator{\Span}{span}

% --------------------------------------------------------------------------
% Custom commands -- linear algebra operators
% --------------------------------------------------------------------------
\DeclareMathOperator{\nullity}{nullity}
\DeclareMathOperator{\range}{range}
\DeclareMathOperator{\nullspace}{null}
\DeclareMathOperator{\proj}{proj}

% --------------------------------------------------------------------------
% Custom commands -- measure theory
% --------------------------------------------------------------------------
\newcommand{\powerset}{\mathcal{P}}
\newcommand{\Borel}{\mathcal{B}}
 (from subdomain directories)

% ---- Encoding and fonts ----
\usepackage[utf8]{inputenc}
\usepackage[T1]{fontenc}

% ---- Mathematics ----
\usepackage{amsmath,amssymb,amsthm}
\usepackage{mathtools}

% ---- Page geometry ----
\usepackage[margin=1in,headheight=14pt]{geometry}

% ---- Hyperlinks ----
\usepackage[colorlinks=true,linkcolor=red,citecolor=blue,urlcolor=blue]{hyperref}

% ---- Lists ----
\usepackage{enumitem}

% ---- Colors and boxes ----
\usepackage{xcolor}
\usepackage[theorems,skins,breakable]{tcolorbox}

% ---- Header/Footer ----
\usepackage{fancyhdr}
\renewcommand{\headrulewidth}{0.4pt}

% ---- Tables ----
\usepackage{booktabs}

% ---- Bibliography ----
\usepackage{natbib}
\bibliographystyle{plainnat}

% --------------------------------------------------------------------------
% Theorem environments
% --------------------------------------------------------------------------
\theoremstyle{plain}
\newtheorem{theorem}{Theorem}[section]
\newtheorem{lemma}[theorem]{Lemma}
\newtheorem{proposition}[theorem]{Proposition}
\newtheorem{corollary}[theorem]{Corollary}

\theoremstyle{definition}
\newtheorem{definition}[theorem]{Definition}
\newtheorem{example}[theorem]{Example}
\newtheorem{exercise}[theorem]{Exercise}
\newtheorem{assumption}[theorem]{Assumption}

\theoremstyle{remark}
\newtheorem{remark}[theorem]{Remark}

% --------------------------------------------------------------------------
% Colored boxes
% --------------------------------------------------------------------------
\newtcolorbox{keyresult}[1][]{%
  colback=green!5!white,
  colframe=green!50!black,
  fonttitle=\bfseries,
  title={Key Result},
  breakable,
  #1
}

\newtcolorbox{rlconnection}[1][]{%
  colback=blue!5!white,
  colframe=blue!50!black,
  fonttitle=\bfseries,
  title={RL/ML Connection},
  breakable,
  #1
}

\newtcolorbox{warning}[1][]{%
  colback=red!5!white,
  colframe=red!50!black,
  fonttitle=\bfseries,
  title={Warning},
  breakable,
  #1
}

% --------------------------------------------------------------------------
% Custom commands -- number sets
% --------------------------------------------------------------------------
\newcommand{\R}{\mathbb{R}}
\newcommand{\N}{\mathbb{N}}
\newcommand{\Q}{\mathbb{Q}}
\newcommand{\Z}{\mathbb{Z}}
\newcommand{\C}{\mathbb{C}}
\newcommand{\F}{\mathbb{F}}

% --------------------------------------------------------------------------
% Custom commands -- probability and expectation
% --------------------------------------------------------------------------
\newcommand{\E}{\mathbb{E}}
\newcommand{\Prob}{\mathbb{P}}
\newcommand{\Var}{\operatorname{Var}}
\newcommand{\Cov}{\operatorname{Cov}}
\newcommand{\indicator}{\mathbf{1}}

% --------------------------------------------------------------------------
% Custom commands -- convergence arrows
% --------------------------------------------------------------------------
\newcommand{\cas}{\xrightarrow{\text{a.s.}}}
\newcommand{\cprob}{\xrightarrow{\Prob}}
\newcommand{\clp}[1]{\xrightarrow{L^{#1}}}
\newcommand{\cdist}{\xrightarrow{d}}

% --------------------------------------------------------------------------
% Custom commands -- common operators
% --------------------------------------------------------------------------
\DeclareMathOperator{\dom}{dom}
\DeclareMathOperator{\epi}{epi}
\DeclareMathOperator{\conv}{conv}
\DeclareMathOperator{\relint}{relint}
\DeclareMathOperator{\aff}{aff}
\DeclareMathOperator{\cl}{cl}
\DeclareMathOperator{\interior}{int}
\DeclareMathOperator{\tr}{tr}
\DeclareMathOperator{\diag}{diag}
\DeclareMathOperator{\rank}{rank}
\DeclareMathOperator{\sgn}{sgn}
\DeclareMathOperator{\supp}{supp}
\DeclareMathOperator{\argmin}{arg\,min}
\DeclareMathOperator{\argmax}{arg\,max}
\DeclareMathOperator{\prox}{prox}
\DeclareMathOperator{\dist}{dist}
\DeclareMathOperator{\diam}{diam}
\DeclareMathOperator{\Span}{span}

% --------------------------------------------------------------------------
% Custom commands -- linear algebra operators
% --------------------------------------------------------------------------
\DeclareMathOperator{\nullity}{nullity}
\DeclareMathOperator{\range}{range}
\DeclareMathOperator{\nullspace}{null}
\DeclareMathOperator{\proj}{proj}

% --------------------------------------------------------------------------
% Custom commands -- measure theory
% --------------------------------------------------------------------------
\newcommand{\powerset}{\mathcal{P}}
\newcommand{\Borel}{\mathcal{B}}
 (from subdomain directories)

% ---- Encoding and fonts ----
\usepackage[utf8]{inputenc}
\usepackage[T1]{fontenc}

% ---- Mathematics ----
\usepackage{amsmath,amssymb,amsthm}
\usepackage{mathtools}

% ---- Page geometry ----
\usepackage[margin=1in,headheight=14pt]{geometry}

% ---- Hyperlinks ----
\usepackage[colorlinks=true,linkcolor=red,citecolor=blue,urlcolor=blue]{hyperref}

% ---- Lists ----
\usepackage{enumitem}

% ---- Colors and boxes ----
\usepackage{xcolor}
\usepackage[theorems,skins,breakable]{tcolorbox}

% ---- Header/Footer ----
\usepackage{fancyhdr}
\renewcommand{\headrulewidth}{0.4pt}

% ---- Tables ----
\usepackage{booktabs}

% ---- Bibliography ----
\usepackage{natbib}
\bibliographystyle{plainnat}

% --------------------------------------------------------------------------
% Theorem environments
% --------------------------------------------------------------------------
\theoremstyle{plain}
\newtheorem{theorem}{Theorem}[section]
\newtheorem{lemma}[theorem]{Lemma}
\newtheorem{proposition}[theorem]{Proposition}
\newtheorem{corollary}[theorem]{Corollary}

\theoremstyle{definition}
\newtheorem{definition}[theorem]{Definition}
\newtheorem{example}[theorem]{Example}
\newtheorem{exercise}[theorem]{Exercise}
\newtheorem{assumption}[theorem]{Assumption}

\theoremstyle{remark}
\newtheorem{remark}[theorem]{Remark}

% --------------------------------------------------------------------------
% Colored boxes
% --------------------------------------------------------------------------
\newtcolorbox{keyresult}[1][]{%
  colback=green!5!white,
  colframe=green!50!black,
  fonttitle=\bfseries,
  title={Key Result},
  breakable,
  #1
}

\newtcolorbox{rlconnection}[1][]{%
  colback=blue!5!white,
  colframe=blue!50!black,
  fonttitle=\bfseries,
  title={RL/ML Connection},
  breakable,
  #1
}

\newtcolorbox{warning}[1][]{%
  colback=red!5!white,
  colframe=red!50!black,
  fonttitle=\bfseries,
  title={Warning},
  breakable,
  #1
}

% --------------------------------------------------------------------------
% Custom commands -- number sets
% --------------------------------------------------------------------------
\newcommand{\R}{\mathbb{R}}
\newcommand{\N}{\mathbb{N}}
\newcommand{\Q}{\mathbb{Q}}
\newcommand{\Z}{\mathbb{Z}}
\newcommand{\C}{\mathbb{C}}
\newcommand{\F}{\mathbb{F}}

% --------------------------------------------------------------------------
% Custom commands -- probability and expectation
% --------------------------------------------------------------------------
\newcommand{\E}{\mathbb{E}}
\newcommand{\Prob}{\mathbb{P}}
\newcommand{\Var}{\operatorname{Var}}
\newcommand{\Cov}{\operatorname{Cov}}
\newcommand{\indicator}{\mathbf{1}}

% --------------------------------------------------------------------------
% Custom commands -- convergence arrows
% --------------------------------------------------------------------------
\newcommand{\cas}{\xrightarrow{\text{a.s.}}}
\newcommand{\cprob}{\xrightarrow{\Prob}}
\newcommand{\clp}[1]{\xrightarrow{L^{#1}}}
\newcommand{\cdist}{\xrightarrow{d}}

% --------------------------------------------------------------------------
% Custom commands -- common operators
% --------------------------------------------------------------------------
\DeclareMathOperator{\dom}{dom}
\DeclareMathOperator{\epi}{epi}
\DeclareMathOperator{\conv}{conv}
\DeclareMathOperator{\relint}{relint}
\DeclareMathOperator{\aff}{aff}
\DeclareMathOperator{\cl}{cl}
\DeclareMathOperator{\interior}{int}
\DeclareMathOperator{\tr}{tr}
\DeclareMathOperator{\diag}{diag}
\DeclareMathOperator{\rank}{rank}
\DeclareMathOperator{\sgn}{sgn}
\DeclareMathOperator{\supp}{supp}
\DeclareMathOperator{\argmin}{arg\,min}
\DeclareMathOperator{\argmax}{arg\,max}
\DeclareMathOperator{\prox}{prox}
\DeclareMathOperator{\dist}{dist}
\DeclareMathOperator{\diam}{diam}
\DeclareMathOperator{\Span}{span}

% --------------------------------------------------------------------------
% Custom commands -- linear algebra operators
% --------------------------------------------------------------------------
\DeclareMathOperator{\nullity}{nullity}
\DeclareMathOperator{\range}{range}
\DeclareMathOperator{\nullspace}{null}
\DeclareMathOperator{\proj}{proj}

% --------------------------------------------------------------------------
% Custom commands -- measure theory
% --------------------------------------------------------------------------
\newcommand{\powerset}{\mathcal{P}}
\newcommand{\Borel}{\mathcal{B}}

\pagestyle{fancy}
\fancyhf{}
\fancyhead[L]{\textsc{Lesson 7: Inner Products and Spectral Theory}}
\fancyhead[R]{\thepage}
\title{Lesson 7: Inner Products and Spectral Theory}
\author{AI/ML Mathematics}
\date{}
\begin{document}
\maketitle
\tableofcontents
\newpage

%% ============================================================
%% SECTION 1: MOTIVATION AND OVERVIEW
%% ============================================================
\section{Motivation and Overview}
\label{sec:motivation}

Suppose we wish to approximate a value function $V^\pi$ in reinforcement
learning using a linear combination of $d$ basis functions
$\phi_1, \ldots, \phi_d$.  Among all vectors in the subspace
$U = \Span(\phi_1, \ldots, \phi_d)$, which one is closest to $V^\pi$?
This question has no answer without a notion of ``distance,'' and distance
requires an inner product.  Once we equip the space with an inner product,
the best approximation is the \emph{orthogonal projection} of $V^\pi$ onto
$U$ -- a result that underpins least-squares methods, principal component
analysis, and the natural policy gradient.

Beyond approximation, inner products unlock the most powerful structural
theorem in finite-dimensional linear algebra: the \emph{spectral theorem}
for real symmetric matrices.  This theorem guarantees that every real
symmetric matrix has an orthonormal basis of eigenvectors with real
eigenvalues -- a fact that is foundational for PCA, kernel methods,
Hessian analysis, and the Fisher information matrix in natural gradient
methods.  The tools developed in this lesson will let us define inner
products, prove the spectral theorem, characterize positive definiteness,
and establish the Rayleigh quotient bounds that govern convergence rates.

\begin{rlconnection}[title={Where Inner Products and Spectral Theory Appear in RL/ML}]
\begin{itemize}[nosep]
  \item \textbf{Least squares and projection.}
        Linear value function approximation finds $\hat{V} = \Phi w$
        minimizing $\|V^\pi - \Phi w\|_2^2$; the solution is the
        orthogonal projection of $V^\pi$ onto $\range(\Phi)$.
  \item \textbf{PCA and spectral decomposition.}
        Principal component analysis diagonalizes the covariance matrix
        $\Sigma = \frac{1}{n}\sum_i x_i x_i^\top$ using the spectral
        theorem to find the directions of maximum variance.
  \item \textbf{Condition number and convergence.}
        The condition number $\kappa(A) = \|A\| \cdot \|A^{-1}\|$
        governs the convergence rate of iterative methods; for symmetric
        positive definite $A$, $\kappa = \lambda_{\max}/\lambda_{\min}$.
  \item \textbf{Natural policy gradient.}
        The Fisher information matrix $F(\theta)$ defines an inner product
        on the parameter space; the natural gradient
        $\tilde{\nabla} J = F^{-1} \nabla J$ is the steepest ascent
        direction in this geometry, computed via spectral decomposition.
\end{itemize}
\end{rlconnection}

\paragraph{Prerequisites.}
Determinants, eigenvalues, characteristic polynomial, algebraic and
geometric multiplicity, diagonalizability criterion, Cayley--Hamilton
theorem (Lesson~6).

\paragraph{Learning objectives.}
After completing this lesson, the student will be able to:
\begin{itemize}[nosep]
  \item \textbf{Define} inner products, norms, and orthogonality, and
        verify the inner product axioms in concrete examples.
  \item \textbf{Prove} the Cauchy--Schwarz inequality and apply the
        Gram--Schmidt process to construct orthonormal bases.
  \item \textbf{Derive} the orthogonal projection formula and prove
        the best approximation theorem.
  \item \textbf{Prove} the spectral theorem for real symmetric matrices:
        every such matrix is orthogonally diagonalizable.
  \item \textbf{Classify} symmetric matrices as positive definite,
        positive semidefinite, or indefinite using eigenvalue criteria.
  \item \textbf{Apply} the Rayleigh quotient and Courant--Fischer theorem
        to characterize extreme eigenvalues.
\end{itemize}

%% ============================================================
%% SECTION 2: REFERENCES
%% ============================================================
\section{References}
\label{sec:references}

\begin{itemize}
  \item \citet[Chapters 6--7]{axler2015} -- Inner product spaces,
        orthonormal bases, the spectral theorem for self-adjoint operators.
  \item \citet[Chapter 8]{hoffman1971} -- Inner product spaces and the
        principal axis theorem with detailed proofs.
  \item \citet[Chapters 6--7]{strang2016} -- Computational perspective
        on orthogonality, Gram--Schmidt, and symmetric matrices.
  \item \citet[Chapter 4]{horn2013} -- Spectral theory for Hermitian
        matrices, Rayleigh quotient, Courant--Fischer, and Weyl
        inequalities.
  \item \citet[Chapter 8]{lax2007} -- Spectral theory with connections
        to functional analysis and operator theory.
\end{itemize}

%% ============================================================
%% SECTION 3: INNER PRODUCT SPACES
%% ============================================================
\section{Inner Product Spaces}
\label{sec:inner_products}

\subsection{Definition and Examples}

\begin{definition}[Inner Product]
\label{def:inner_product}
An \emph{inner product} on a real vector space $V$ is a function
$\langle \cdot, \cdot \rangle\colon V \times V \to \R$ satisfying, for
all $u, v, w \in V$ and $\alpha \in \R$:
\begin{enumerate}[nosep]
  \item \textbf{Symmetry:} $\langle u, v \rangle = \langle v, u \rangle$.
  \item \textbf{Linearity in the first argument:}
        $\langle \alpha u + v, w \rangle
        = \alpha \langle u, w \rangle + \langle v, w \rangle$.
  \item \textbf{Positive definiteness:}
        $\langle v, v \rangle \geq 0$, with equality if and only if $v = 0$.
\end{enumerate}
A vector space equipped with an inner product is called an
\emph{inner product space}.
\end{definition}

\begin{remark}
\label{rem:bilinearity}
Symmetry and linearity in the first argument together imply linearity in
the second argument (bilinearity):
$\langle u, \alpha v + w \rangle = \langle \alpha v + w, u \rangle
= \alpha \langle v, u \rangle + \langle w, u \rangle
= \alpha \langle u, v \rangle + \langle u, w \rangle$.
\end{remark}

\begin{example}[Standard Inner Products]
\label{ex:standard_ips}
\begin{enumerate}[nosep]
  \item \textbf{Euclidean inner product on $\R^n$:}
        $\langle x, y \rangle = \sum_{i=1}^n x_i y_i = x^\top y$.
  \item \textbf{Weighted inner product on $\R^n$:}
        For a symmetric positive definite matrix $M \in \R^{n \times n}$,
        $\langle x, y \rangle_M = x^\top M y$ defines an inner product.
  \item \textbf{$L^2$ inner product on $C[a,b]$:}
        $\langle f, g \rangle = \int_a^b f(t) g(t)\, dt$ for continuous
        functions $f, g\colon [a,b] \to \R$.
\end{enumerate}
\end{example}

\begin{definition}[Norm Induced by an Inner Product]
\label{def:norm}
The \emph{norm} induced by an inner product is
$\|v\| = \sqrt{\langle v, v \rangle}$.  The \emph{distance} between
$u$ and $v$ is $\|u - v\|$.
\end{definition}

\subsection{Cauchy--Schwarz and Triangle Inequalities}

\begin{theorem}[Cauchy--Schwarz Inequality]
\label{thm:cauchy_schwarz}
For all $u, v$ in an inner product space,
\[
  |\langle u, v \rangle| \leq \|u\| \cdot \|v\|.
\]
Equality holds if and only if $u$ and $v$ are linearly dependent.
\end{theorem}

\begin{proof}
If $v = 0$, both sides are zero.  Assume $v \neq 0$.  For any $t \in \R$,
\[
  0 \leq \|u - tv\|^2
  = \langle u - tv, u - tv \rangle
  = \|u\|^2 - 2t\langle u, v \rangle + t^2 \|v\|^2.
\]
This is a quadratic in $t$ that is nonneg for all $t$.  Minimizing over
$t$ at $t^* = \langle u, v \rangle / \|v\|^2$ gives
\[
  0 \leq \|u\|^2 - \frac{\langle u, v \rangle^2}{\|v\|^2},
\]
which rearranges to $\langle u, v \rangle^2 \leq \|u\|^2 \|v\|^2$.
Taking square roots yields $|\langle u, v \rangle| \leq \|u\| \cdot \|v\|$.

Equality holds iff $\|u - t^* v\| = 0$, i.e., $u = t^* v$ (linear dependence).
\end{proof}

\begin{corollary}[Triangle Inequality]
\label{cor:triangle}
For all $u, v$ in an inner product space, $\|u + v\| \leq \|u\| + \|v\|$.
\end{corollary}

\begin{proof}
\begin{align*}
  \|u + v\|^2
  &= \|u\|^2 + 2\langle u, v \rangle + \|v\|^2
   \leq \|u\|^2 + 2\|u\|\,\|v\| + \|v\|^2
   = (\|u\| + \|v\|)^2. \qedhere
\end{align*}
\end{proof}

With the Cauchy--Schwarz inequality established, we can develop the
geometric structure of inner product spaces: orthogonality, projections,
and orthonormal bases.

%% ============================================================
%% SECTION 4: ORTHOGONALITY AND ORTHONORMAL BASES
%% ============================================================
\section{Orthogonality and Orthonormal Bases}
\label{sec:orthogonality}

\subsection{Orthogonal and Orthonormal Lists}

\begin{definition}[Orthogonality]
\label{def:orthogonal}
Two vectors $u, v$ in an inner product space are \emph{orthogonal},
written $u \perp v$, if $\langle u, v \rangle = 0$.  A list of vectors
$v_1, \ldots, v_m$ is \emph{orthogonal} if $\langle v_i, v_j \rangle = 0$
for all $i \neq j$, and \emph{orthonormal} if additionally
$\|v_i\| = 1$ for each $i$.
\end{definition}

\begin{theorem}[Pythagorean Theorem]
\label{thm:pythagoras}
If $v_1, \ldots, v_m$ are pairwise orthogonal, then
\[
  \|v_1 + \cdots + v_m\|^2 = \|v_1\|^2 + \cdots + \|v_m\|^2.
\]
\end{theorem}

\begin{proof}
$\|v_1 + \cdots + v_m\|^2 = \langle \sum_i v_i, \sum_j v_j \rangle
= \sum_i \sum_j \langle v_i, v_j \rangle
= \sum_i \|v_i\|^2$,
since all cross terms vanish by orthogonality.
\end{proof}

\begin{proposition}[Orthogonal Lists Are Independent]
\label{prop:orth_indep}
Any orthogonal list of nonzero vectors is linearly independent.
\end{proposition}

\begin{proof}
Let $v_1, \ldots, v_m$ be orthogonal and nonzero.  If
$c_1 v_1 + \cdots + c_m v_m = 0$, take the inner product with $v_k$:
$c_k \|v_k\|^2 = 0$.  Since $\|v_k\| \neq 0$, $c_k = 0$ for each $k$.
\end{proof}

\begin{example}[Orthonormal Basis of $\R^3$]
\label{ex:onb_R3}
The standard basis $e_1 = (1,0,0)$, $e_2 = (0,1,0)$, $e_3 = (0,0,1)$
is orthonormal with respect to the Euclidean inner product.  For any
$v = (a,b,c) \in \R^3$, the coordinates are recovered by inner products:
$a = \langle v, e_1 \rangle$, $b = \langle v, e_2 \rangle$,
$c = \langle v, e_3 \rangle$.
\end{example}

\subsection{The Gram--Schmidt Process}

\begin{theorem}[Gram--Schmidt Orthogonalization]
\label{thm:gram_schmidt}
Let $v_1, \ldots, v_m$ be a linearly independent list in an inner product
space $V$.  Then there exists an orthonormal list $e_1, \ldots, e_m$ such
that
\[
  \Span(e_1, \ldots, e_k) = \Span(v_1, \ldots, v_k)
  \quad \text{for each } k = 1, \ldots, m.
\]
The orthonormal list is constructed by:
\begin{align*}
  w_1 &= v_1, \quad e_1 = w_1 / \|w_1\|, \\
  w_k &= v_k - \sum_{j=1}^{k-1} \langle v_k, e_j \rangle\, e_j,
         \quad e_k = w_k / \|w_k\|,
         \quad k = 2, \ldots, m.
\end{align*}
\end{theorem}

\begin{proof}
We proceed by induction.  \textbf{Base case:} $v_1 \neq 0$ (by
independence), so $w_1 = v_1 \neq 0$ and $e_1 = v_1/\|v_1\|$ is a
unit vector.  Clearly $\Span(e_1) = \Span(v_1)$.

\textbf{Inductive step:} Assume $e_1, \ldots, e_{k-1}$ is orthonormal
with $\Span(e_1, \ldots, e_{j}) = \Span(v_1, \ldots, v_j)$ for
$j \leq k-1$.  Define
$w_k = v_k - \sum_{j=1}^{k-1} \langle v_k, e_j \rangle\, e_j$.
For each $i \leq k-1$:
\[
  \langle w_k, e_i \rangle
  = \langle v_k, e_i \rangle
    - \sum_{j=1}^{k-1} \langle v_k, e_j \rangle \langle e_j, e_i \rangle
  = \langle v_k, e_i \rangle - \langle v_k, e_i \rangle = 0.
\]
Also $w_k \neq 0$: if $w_k = 0$, then
$v_k = \sum_j \langle v_k, e_j \rangle\, e_j \in \Span(e_1, \ldots, e_{k-1})
= \Span(v_1, \ldots, v_{k-1})$, contradicting the independence of the
original list.  Set $e_k = w_k / \|w_k\|$.  Then $e_1, \ldots, e_k$ is
orthonormal, and $\Span(e_1, \ldots, e_k) = \Span(v_1, \ldots, v_k)$
(since $w_k$ is a combination of $v_k$ and $e_1, \ldots, e_{k-1}$, and
$v_k$ can be recovered from $w_k$ and $e_1, \ldots, e_{k-1}$).
\end{proof}

\begin{example}[Gram--Schmidt in $\R^3$]
\label{ex:gram_schmidt_R3}
Apply Gram--Schmidt to $v_1 = (1, 1, 0)$, $v_2 = (1, 0, 1)$.

\textbf{Step 1:} $w_1 = (1,1,0)$, $\|w_1\| = \sqrt{2}$,
$e_1 = \frac{1}{\sqrt{2}}(1, 1, 0)$.

\textbf{Step 2:} $\langle v_2, e_1 \rangle = 1/\sqrt{2}$.
$w_2 = (1,0,1) - \frac{1}{2}(1,1,0) = (1/2, -1/2, 1)$.
$\|w_2\| = \sqrt{3/2}$,
$e_2 = \frac{1}{\sqrt{6}}(1, -1, 2)$.
One verifies $\langle e_1, e_2 \rangle = 0$ and
$\Span(e_1, e_2) = \Span(v_1, v_2)$.
\end{example}

\begin{corollary}
\label{cor:onb_exists}
Every finite-dimensional inner product space has an orthonormal basis.
\end{corollary}

\begin{proof}
Start with any basis and apply the Gram--Schmidt process
(Theorem~\ref{thm:gram_schmidt}).
\end{proof}

\subsection{Orthogonal Complement and Decomposition}

\begin{definition}[Orthogonal Complement]
\label{def:orth_complement}
If $U$ is a subspace of an inner product space $V$, the \emph{orthogonal
complement} of $U$ is
\[
  U^\perp = \{v \in V : \langle v, u \rangle = 0 \text{ for all } u \in U\}.
\]
\end{definition}

\begin{proposition}
\label{prop:orth_complement_subspace}
$U^\perp$ is a subspace of $V$, and $U \cap U^\perp = \{0\}$.
\end{proposition}

\begin{proof}
$U^\perp$ is a subspace: $0 \in U^\perp$, and closure under linear
combinations follows from bilinearity.  If $v \in U \cap U^\perp$, then
$\langle v, v \rangle = 0$, so $v = 0$.
\end{proof}

\begin{theorem}[Orthogonal Decomposition]
\label{thm:orth_decomp}
If $V$ is a finite-dimensional inner product space and $U$ is a subspace,
then $V = U \oplus U^\perp$.  Equivalently, every $v \in V$ can be
written uniquely as $v = u + w$ with $u \in U$ and $w \in U^\perp$.
\end{theorem}

\begin{proof}
Let $e_1, \ldots, e_k$ be an orthonormal basis of $U$ (exists by
Corollary~\ref{cor:onb_exists}).  For any $v \in V$, set
\[
  u = \sum_{j=1}^{k} \langle v, e_j \rangle\, e_j, \qquad
  w = v - u.
\]
Then $u \in U$ and for each $i = 1, \ldots, k$:
\[
  \langle w, e_i \rangle
  = \langle v, e_i \rangle - \sum_{j=1}^{k} \langle v, e_j \rangle
    \langle e_j, e_i \rangle
  = \langle v, e_i \rangle - \langle v, e_i \rangle = 0.
\]
Hence $w \perp e_i$ for each basis element of $U$, so $w \in U^\perp$.
This shows $V = U + U^\perp$.  By Proposition~\ref{prop:orth_complement_subspace},
$U \cap U^\perp = \{0\}$, so the sum is direct.
\end{proof}

\begin{corollary}[Dimension Formula]
\label{cor:dim_orth}
$\dim U + \dim U^\perp = \dim V$ and $(U^\perp)^\perp = U$.
\end{corollary}

\begin{proof}
The first follows from $V = U \oplus U^\perp$.  For the second,
$U \subseteq (U^\perp)^\perp$ is clear, and both have dimension $\dim U$.
\end{proof}

\subsection{Orthogonal Projection and Best Approximation}

\begin{definition}[Orthogonal Projection]
\label{def:projection}
Let $U$ be a subspace of a finite-dimensional inner product space $V$.
The \emph{orthogonal projection} onto $U$ is the map
$\proj_U\colon V \to V$ defined by $\proj_U(v) = u$, where $v = u + w$
is the orthogonal decomposition with $u \in U$ and $w \in U^\perp$.
If $e_1, \ldots, e_k$ is an orthonormal basis of $U$, then
\[
  \proj_U(v) = \sum_{j=1}^{k} \langle v, e_j \rangle\, e_j.
\]
\end{definition}

\begin{keyresult}[title={Orthogonal Projection and Best Approximation}]
\begin{theorem}[Best Approximation]
\label{thm:best_approx}
Let $U$ be a subspace of a finite-dimensional inner product space $V$
and let $v \in V$.  Then $\proj_U(v)$ is the unique element of $U$ that
minimizes the distance to $v$:
\[
  \|v - \proj_U(v)\| < \|v - u\| \quad \text{for all } u \in U,
  \; u \neq \proj_U(v).
\]
\end{theorem}
\end{keyresult}

\begin{proof}
Let $\hat{u} = \proj_U(v)$, so $v - \hat{u} \in U^\perp$.  For any
$u \in U$:
\[
  \|v - u\|^2 = \|(v - \hat{u}) + (\hat{u} - u)\|^2
  = \|v - \hat{u}\|^2 + \|\hat{u} - u\|^2,
\]
where the cross term vanishes because $v - \hat{u} \in U^\perp$ and
$\hat{u} - u \in U$ (Pythagorean theorem, Theorem~\ref{thm:pythagoras}).
Hence $\|v - u\|^2 \geq \|v - \hat{u}\|^2$, with equality iff
$u = \hat{u}$.
\end{proof}

\begin{example}[Projection onto a Line]
\label{ex:proj_line}
Let $U = \Span(a)$ for $a \neq 0$ in $\R^n$.  The orthonormal basis of
$U$ is $e_1 = a/\|a\|$, so
\[
  \proj_U(v) = \frac{\langle v, a \rangle}{\|a\|^2}\, a
  = \frac{a^\top v}{a^\top a}\, a.
\]
For $a = (1,1)^\top$ and $v = (3,1)^\top$ in $\R^2$:
$\proj_U(v) = \frac{4}{2}(1,1)^\top = (2,2)^\top$.
The residual $v - \proj_U(v) = (1,-1)^\top$ is orthogonal to $a$:
$(1,-1) \cdot (1,1) = 0$.
\end{example}

\begin{rlconnection}[title={Least Squares and Linear Value Function Approximation}]
In linear value function approximation, states are represented by feature
vectors $\phi(s_i) \in \R^d$, forming a matrix $\Phi \in \R^{n \times d}$.
The projected Bellman equation seeks $w^* \in \R^d$ such that
$\Phi w^*$ is the orthogonal projection of $r + \gamma P V$ onto
$\range(\Phi)$ under the inner product
$\langle u, v \rangle_\mu = \sum_i \mu(s_i) u_i v_i$ weighted by the
stationary distribution $\mu$.  By the best approximation theorem
(Theorem~\ref{thm:best_approx}), this projection minimizes
$\|V - \Phi w\|_\mu^2$ over all $w \in \R^d$.  The normal equations
$\Phi^\top D_\mu \Phi\, w = \Phi^\top D_\mu (r + \gamma P V)$ are the
concrete form of $v - \proj_U(v) \in U^\perp$.
\end{rlconnection}

\subsection{Bessel's Inequality}

\begin{theorem}[Bessel's Inequality]
\label{thm:bessel}
Let $e_1, \ldots, e_m$ be an orthonormal list (not necessarily a basis)
in an inner product space $V$.  For any $v \in V$,
\[
  \sum_{j=1}^{m} |\langle v, e_j \rangle|^2 \leq \|v\|^2.
\]
Equality holds if and only if $v \in \Span(e_1, \ldots, e_m)$.
\end{theorem}

\begin{proof}
With $U = \Span(e_1, \ldots, e_m)$, the Pythagorean theorem gives
$\|v\|^2 = \|\proj_U(v)\|^2 + \|v - \proj_U(v)\|^2
= \sum_j |\langle v, e_j \rangle|^2 + \|v - \proj_U(v)\|^2
\geq \sum_j |\langle v, e_j \rangle|^2$,
with equality iff $v \in U$.
\end{proof}

%% ============================================================
%% SECTION 5: THE SPECTRAL THEOREM
%% ============================================================
\section{The Spectral Theorem for Real Symmetric Matrices}
\label{sec:spectral}

We now turn to the central result of this lesson.  Throughout this
section, all matrices are real ($\F = \R$) and the inner product is
the standard Euclidean inner product $\langle x, y \rangle = x^\top y$.

\subsection{Eigenvalues of Symmetric Matrices Are Real}

\begin{lemma}
\label{lem:symmetric_real_eigen}
Every eigenvalue of a real symmetric matrix $A \in \R^{n \times n}$ is
real.
\end{lemma}

\begin{proof}
We work over $\C$ temporarily.  Let $\lambda \in \C$ be an eigenvalue
with eigenvector $v \in \C^n \setminus \{0\}$: $Av = \lambda v$.  Take
the conjugate transpose: $\bar{v}^\top A^\top = \bar{\lambda} \bar{v}^\top$.
Since $A = A^\top$ (and $A$ has real entries, so $\bar{A} = A$):
\[
  \bar{v}^\top A v = \bar{\lambda}\, \bar{v}^\top v.
\]
Also, $\bar{v}^\top A v = \bar{v}^\top (\lambda v) = \lambda\, \bar{v}^\top v$.
Since $\bar{v}^\top v = \sum_i |v_i|^2 > 0$, we can divide to get
$\lambda = \bar{\lambda}$, so $\lambda \in \R$.
\end{proof}

\subsection{Orthogonality of Eigenvectors}

\begin{lemma}
\label{lem:symmetric_orth_eigen}
Let $A \in \R^{n \times n}$ be symmetric.  Eigenvectors corresponding to
distinct eigenvalues are orthogonal.
\end{lemma}

\begin{proof}
Let $Av_1 = \lambda_1 v_1$ and $Av_2 = \lambda_2 v_2$ with
$\lambda_1 \neq \lambda_2$.  Then
\[
  \lambda_1 \langle v_1, v_2 \rangle
  = \langle Av_1, v_2 \rangle
  = v_1^\top A^\top v_2
  = v_1^\top A v_2
  = \langle v_1, Av_2 \rangle
  = \lambda_2 \langle v_1, v_2 \rangle.
\]
Since $\lambda_1 \neq \lambda_2$, we must have
$\langle v_1, v_2 \rangle = 0$.
\end{proof}

\subsection{Invariant Subspaces and the Spectral Theorem}

\begin{lemma}[Symmetric Matrices Preserve Orthogonal Complements]
\label{lem:symmetric_invariant}
Let $A \in \R^{n \times n}$ be symmetric and let $U$ be an
$A$-invariant subspace (i.e., $Au \in U$ for all $u \in U$).  Then
$U^\perp$ is also $A$-invariant.
\end{lemma}

\begin{proof}
Let $w \in U^\perp$.  For any $u \in U$, we have $Au \in U$ (by
$A$-invariance) and
\[
  \langle Aw, u \rangle = \langle w, A^\top u \rangle
  = \langle w, Au \rangle = 0,
\]
since $Au \in U$ and $w \in U^\perp$.  Hence $Aw \perp U$, i.e.,
$Aw \in U^\perp$.
\end{proof}

\begin{keyresult}[title={The Spectral Theorem for Real Symmetric Matrices}]
\begin{theorem}[Spectral Theorem]
\label{thm:spectral}
Let $A \in \R^{n \times n}$ be symmetric.  Then:
\begin{enumerate}[nosep]
  \item All eigenvalues of $A$ are real.
  \item There exists an orthonormal basis of $\R^n$ consisting of
        eigenvectors of $A$.
  \item There exists an orthogonal matrix $Q$ (i.e., $Q^\top Q = I$)
        such that $Q^\top A Q = \diag(\lambda_1, \ldots, \lambda_n)$.
\end{enumerate}
Equivalently, $A$ admits the \emph{spectral decomposition}
$A = \sum_{i=1}^{n} \lambda_i\, q_i q_i^\top$, where
$q_1, \ldots, q_n$ are orthonormal eigenvectors.
\end{theorem}
\end{keyresult}

\begin{proof}
(1) is Lemma~\ref{lem:symmetric_real_eigen}.

(2) We prove by induction on $n$.  For $n = 1$, any nonzero vector is
an eigenvector, and normalizing gives an orthonormal basis.

For the inductive step, let $n \geq 2$.  Since the characteristic
polynomial $p_A$ has degree $n$ and real coefficients, and all roots
are real (by part (1)), $A$ has at least one real eigenvalue $\lambda_1$
with unit eigenvector $q_1$.  Let $U = \Span(q_1)$.  Since
$AU = \lambda_1 U \subseteq U$, the subspace $U$ is $A$-invariant.
By Lemma~\ref{lem:symmetric_invariant}, $U^\perp$ is also $A$-invariant.

Let $B = A|_{U^\perp}$ denote the restriction of $A$ to $U^\perp$.
Since $U^\perp$ has dimension $n - 1$ and $B$ is symmetric (as the
restriction of a symmetric operator to an invariant subspace), the
inductive hypothesis provides an orthonormal basis
$q_2, \ldots, q_n$ of $U^\perp$ consisting of eigenvectors of $B$
(and hence of $A$).  Then $q_1, q_2, \ldots, q_n$ is an orthonormal
basis of $\R^n = U \oplus U^\perp$ consisting of eigenvectors of $A$.

(3) Setting $Q = (q_1 \mid q_2 \mid \cdots \mid q_n)$, the columns are
orthonormal, so $Q^\top Q = I$.  The eigenvector equation
$AQ = Q\diag(\lambda_1, \ldots, \lambda_n)$ gives
$Q^\top A Q = \diag(\lambda_1, \ldots, \lambda_n)$.

The spectral decomposition follows from $A = Q\Lambda Q^\top$ where
$\Lambda = \diag(\lambda_i)$:
$A = \sum_{i=1}^n \lambda_i\, q_i q_i^\top$.
\end{proof}

\begin{example}[Spectral Decomposition of a $2 \times 2$ Matrix]
\label{ex:spectral_2x2}
Let $A = \begin{pmatrix} 3 & 1 \\ 1 & 3 \end{pmatrix}$.  The
characteristic polynomial is $\lambda^2 - 6\lambda + 8 = (\lambda-2)(\lambda-4)$.
Eigenvalues: $\lambda_1 = 2$, $\lambda_2 = 4$.

For $\lambda_1 = 2$: $q_1 = \frac{1}{\sqrt{2}}(1, -1)^\top$.
For $\lambda_2 = 4$: $q_2 = \frac{1}{\sqrt{2}}(1, 1)^\top$.

The spectral decomposition is
$A = 2 q_1 q_1^\top + 4 q_2 q_2^\top
= 2 \cdot \frac{1}{2}\begin{pmatrix} 1 & -1 \\ -1 & 1 \end{pmatrix}
+ 4 \cdot \frac{1}{2}\begin{pmatrix} 1 & 1 \\ 1 & 1 \end{pmatrix}
= \begin{pmatrix} 3 & 1 \\ 1 & 3 \end{pmatrix}$.
\end{example}

\begin{example}[Spectral Decomposition of a $3 \times 3$ Matrix]
\label{ex:spectral_3x3}
Let $A = \begin{pmatrix} 2 & 1 & 0 \\ 1 & 2 & 1 \\ 0 & 1 & 2 \end{pmatrix}$.
Expanding $\det(A - \lambda I) = (2-\lambda)[(2-\lambda)^2 - 2]$ gives
eigenvalues $\lambda_1 = 2 - \sqrt{2}$, $\lambda_2 = 2$,
$\lambda_3 = 2 + \sqrt{2}$ -- all real and distinct, confirming the
spectral theorem.
\end{example}

\begin{rlconnection}[title={PCA and Spectral Decomposition}]
In principal component analysis, given centered data
$x_1, \ldots, x_m \in \R^n$, the sample covariance matrix is
$S = \frac{1}{m}\sum_{i=1}^m x_i x_i^\top$.  Since $S$ is symmetric
and positive semidefinite, the spectral theorem gives
$S = \sum_{j=1}^n \lambda_j q_j q_j^\top$ with $\lambda_j \geq 0$.
The top $k$ principal components are the eigenvectors
$q_1, \ldots, q_k$ (ordered by $\lambda_1 \geq \cdots \geq \lambda_n$).
The projection $\proj_U(x) = \sum_{j=1}^k \langle x, q_j \rangle q_j$
onto $U = \Span(q_1, \ldots, q_k)$ is the best rank-$k$ approximation
of the data by the Eckart--Young theorem.
\end{rlconnection}

%% ============================================================
%% SECTION 6: POSITIVE DEFINITENESS
%% ============================================================
\section{Positive Definiteness and Cholesky Factorization}
\label{sec:positive_definite}

\subsection{Characterizations of Positive Definiteness}

\begin{definition}[Positive Definite and Positive Semidefinite]
\label{def:pd}
A symmetric matrix $A \in \R^{n \times n}$ is:
\begin{itemize}[nosep]
  \item \emph{Positive definite} (written $A \succ 0$) if
        $x^\top A x > 0$ for all nonzero $x \in \R^n$.
  \item \emph{Positive semidefinite} (written $A \succeq 0$) if
        $x^\top A x \geq 0$ for all $x \in \R^n$.
\end{itemize}
\end{definition}

\begin{theorem}[Characterizations of Positive Definiteness]
\label{thm:pd_char}
For a symmetric matrix $A \in \R^{n \times n}$, the following are
equivalent:
\begin{enumerate}[nosep]
  \item $A \succ 0$ (positive definite).
  \item All eigenvalues of $A$ are strictly positive.
  \item $A = R^\top R$ for some invertible matrix $R$.
  \item All leading principal minors of $A$ are strictly positive.
\end{enumerate}
\end{theorem}

\begin{proof}
$(1) \Rightarrow (2)$: If $Av = \lambda v$ with $v \neq 0$, then
$0 < v^\top Av = \lambda \|v\|^2$, so $\lambda > 0$.

$(2) \Rightarrow (3)$: By the spectral theorem, $A = Q\Lambda Q^\top$
with $\lambda_i > 0$.  Set $R = \Lambda^{1/2}Q^\top$; then $R^\top R = A$
and $R$ is invertible.

$(3) \Rightarrow (1)$: $x^\top Ax = \|Rx\|^2 > 0$ for $x \neq 0$ since
$R$ is invertible.

$(1) \Leftrightarrow (4)$: See \citet[Theorem~7.2.5]{horn2013}; the key
idea is that principal submatrices of positive definite matrices are
positive definite.
\end{proof}

\begin{warning}[title={Positive Definiteness Requires Symmetry}]
The non-symmetric matrix $A = \bigl(\begin{smallmatrix} 1 & 2 \\ 0 & 1 \end{smallmatrix}\bigr)$
has eigenvalues $\{1,1\}$ (positive), yet
$x^\top Ax = (x_1 + x_2)^2 = 0$ at $x = (1,-1)^\top$.  The eigenvalue
characterization in Theorem~\ref{thm:pd_char} requires symmetry.
\end{warning}

\subsection{Cholesky Factorization}

\begin{theorem}[Cholesky Factorization]
\label{thm:cholesky}
If $A \in \R^{n \times n}$ is symmetric positive definite, then there
exists a unique lower triangular matrix $L$ with positive diagonal entries
such that $A = LL^\top$.
\end{theorem}

\begin{proof}
\textbf{Existence} (by induction on $n$).  For $n = 1$, $L = (\sqrt{a_{11}})$.
For $n \geq 2$, partition $A$ as
$A = \bigl(\begin{smallmatrix} a_{11} & b^\top \\ b & C \end{smallmatrix}\bigr)$
with $a_{11} > 0$.  Set $\ell_{11} = \sqrt{a_{11}}$ and
$\ell = b/\ell_{11}$.  The Schur complement $S = C - \ell\ell^\top$ is
symmetric positive definite, so by induction $S = L_1 L_1^\top$.  Then
$L = \bigl(\begin{smallmatrix} \ell_{11} & 0 \\ \ell & L_1 \end{smallmatrix}\bigr)$
satisfies $A = LL^\top$.

\textbf{Uniqueness:} If $LL^\top = MM^\top$, then $Q = M^{-1}L$ satisfies
$QQ^\top = I$.  Since $Q$ is lower triangular with positive diagonal
(as a product of lower triangular matrices with positive diagonals), and
orthogonal lower triangular matrices with positive diagonal must equal $I$,
we get $L = M$.
\end{proof}

\begin{example}[Cholesky Factorization]
\label{ex:cholesky}
Let $A = \begin{pmatrix} 4 & 2 \\ 2 & 5 \end{pmatrix}$.  Since
$\det A = 16 > 0$ and $a_{11} = 4 > 0$, $A$ is positive definite.
$\ell_{11} = 2$, $\ell_{21} = 2/2 = 1$, $s_{22} = 5 - 1 = 4$,
$\ell_{22} = 2$.  So
$L = \begin{pmatrix} 2 & 0 \\ 1 & 2 \end{pmatrix}$ and
$LL^\top = \begin{pmatrix} 4 & 2 \\ 2 & 5 \end{pmatrix} = A$.
\end{example}

\begin{rlconnection}[title={Condition Number and Convergence of Iterative Methods}]
For a symmetric positive definite system $Ax = b$, the condition number
$\kappa(A) = \lambda_{\max}(A)/\lambda_{\min}(A)$ determines the rate of
convergence of gradient descent and conjugate gradient methods.
Preconditioning replaces $Ax = b$ with $L^{-1}AL^{-\top}(L^\top x) = L^{-1}b$
using the Cholesky factor $A \approx LL^\top$ of a preconditioner.
A well-chosen preconditioner reduces $\kappa$ and accelerates convergence.
In RL, the Fisher information matrix $F(\theta)$ is symmetric positive
(semi)definite; its condition number affects the convergence of natural
policy gradient methods.
\end{rlconnection}

%% ============================================================
%% SECTION 7: RAYLEIGH QUOTIENT AND COURANT-FISCHER
%% ============================================================
\section{Rayleigh Quotient and Courant--Fischer}
\label{sec:rayleigh}

\subsection{The Rayleigh Quotient}

\begin{definition}[Rayleigh Quotient]
\label{def:rayleigh}
For a symmetric matrix $A \in \R^{n \times n}$, the \emph{Rayleigh
quotient} is the function $R_A\colon \R^n \setminus \{0\} \to \R$ defined
by
\[
  R_A(x) = \frac{x^\top A x}{x^\top x} = \frac{\langle Ax, x \rangle}{\|x\|^2}.
\]
\end{definition}

\begin{theorem}[Extreme Values of the Rayleigh Quotient]
\label{thm:rayleigh_bounds}
Let $A \in \R^{n \times n}$ be symmetric with eigenvalues
$\lambda_1 \leq \lambda_2 \leq \cdots \leq \lambda_n$.  Then
\[
  \lambda_1 = \min_{x \neq 0} R_A(x), \qquad
  \lambda_n = \max_{x \neq 0} R_A(x).
\]
The minimum is attained at any eigenvector for $\lambda_1$, and the
maximum at any eigenvector for $\lambda_n$.
\end{theorem}

\begin{proof}
Substituting $y = Q^\top x$ in the spectral decomposition $A = Q\Lambda Q^\top$:
$R_A(x) = \sum_i \lambda_i y_i^2 / \sum_i y_i^2$,
a weighted average of eigenvalues with nonneg weights summing to 1.
Hence $\lambda_1 \leq R_A(x) \leq \lambda_n$.  Equality is attained at
the corresponding eigenvectors: $R_A(q_1) = \lambda_1$ and
$R_A(q_n) = \lambda_n$.
\end{proof}

\subsection{The Courant--Fischer Minimax Theorem}

\begin{theorem}[Courant--Fischer Minimax Theorem]
\label{thm:courant_fischer}
Let $A \in \R^{n \times n}$ be symmetric with eigenvalues
$\lambda_1 \leq \cdots \leq \lambda_n$.  For each $k = 1, \ldots, n$,
\[
  \lambda_k = \min_{\substack{S \subseteq \R^n \\ \dim S = k}}
              \max_{\substack{x \in S \\ x \neq 0}} R_A(x)
            = \max_{\substack{S \subseteq \R^n \\ \dim S = n-k+1}}
              \min_{\substack{x \in S \\ x \neq 0}} R_A(x).
\]
\end{theorem}

\begin{proof}
We prove $\lambda_k = \min_{\dim S = k} \max_{x \in S \setminus \{0\}} R_A(x)$.
Let $q_1, \ldots, q_n$ be orthonormal eigenvectors with $Aq_i = \lambda_i q_i$.

\textbf{Upper bound:} $S_0 = \Span(q_1, \ldots, q_k)$ has dimension $k$.
For $x = \sum_{i=1}^k c_i q_i \neq 0$,
$R_A(x) = \sum_i \lambda_i c_i^2 / \sum_i c_i^2 \leq \lambda_k$, with
equality at $x = q_k$.

\textbf{Lower bound:} For any $k$-dimensional $S$, let
$T = \Span(q_k, \ldots, q_n)$ with $\dim T = n - k + 1$.  Since
$\dim S + \dim T > n$, there exists $x \in S \cap T \setminus \{0\}$.
Writing $x = \sum_{i=k}^n c_i q_i$ gives $R_A(x) \geq \lambda_k$.

Combining: $\lambda_k = \min_{\dim S = k} \max_{x \in S} R_A(x)$.
The max-min formula is proved analogously.
\end{proof}

\begin{example}[Rayleigh Quotient Bounds]
\label{ex:rayleigh_bound}
For $A = \begin{pmatrix} 4 & 1 \\ 1 & 2 \end{pmatrix}$, evaluating
$R_A$ at $x = (0,1)^\top$ gives $R_A = 2$, so $\lambda_{\min} \leq 2$.
At $x = (1,1)^\top$, $R_A = 4$, so $\lambda_{\max} \geq 4$.  The exact
eigenvalues $3 \pm \sqrt{2} \approx 1.59, 4.41$ are consistent.
\end{example}

\begin{rlconnection}[title={Natural Policy Gradient and the Fisher Matrix}]
In natural policy gradient methods, the update
$\theta_{t+1} = \theta_t + \alpha F(\theta_t)^{-1} \nabla J(\theta_t)$
uses the Fisher information matrix $F(\theta) \succeq 0$.  Its spectral
decomposition enables efficient computation of $F^{-1}\nabla J$, and
the Courant--Fischer theorem characterizes how the curvature of the
KL-divergence landscape varies across parameter subspaces.
\end{rlconnection}

%% ============================================================
%% SECTION 8: EXERCISES
%% ============================================================
\section{Exercises}
\label{sec:exercises}

\begin{exercise}[Inner Product Verification]
\label{exer:ip_verify}
\textnormal{(\(*\))} \\
Let $V = \R^2$ and define $\langle x, y \rangle = 2x_1 y_1 + x_1 y_2 + x_2 y_1 + 3x_2 y_2$.
\begin{enumerate}[nosep,label=(\alph*)]
  \item Verify that this is an inner product by checking all axioms.
  \item Find the matrix $M$ such that $\langle x, y \rangle = x^\top M y$.
  \item Compute $\|e_1\|$ and $\|e_2\|$ in this inner product.
\end{enumerate}
\end{exercise}

\begin{exercise}[Gram--Schmidt and Projection]
\label{exer:gram_schmidt}
\textnormal{(\(*\))} \\
In $\R^3$ with the standard inner product:
\begin{enumerate}[nosep,label=(\alph*)]
  \item Apply the Gram--Schmidt process to $v_1 = (1, 1, 1)$,
        $v_2 = (1, 1, 0)$, $v_3 = (1, 0, 0)$ to produce an orthonormal
        basis.
  \item Compute the orthogonal projection of $w = (2, 3, 5)$ onto
        $U = \Span(v_1, v_2)$.
  \item Verify that $w - \proj_U(w)$ is orthogonal to both $v_1$ and $v_2$.
\end{enumerate}
\end{exercise}

\begin{exercise}[Spectral Decomposition]
\label{exer:spectral}
\textnormal{(\(**\))} \\
Let $A = \begin{pmatrix} 5 & 2 \\ 2 & 2 \end{pmatrix}$.
\begin{enumerate}[nosep,label=(\alph*)]
  \item Find all eigenvalues and a corresponding orthonormal eigenbasis.
  \item Write the spectral decomposition $A = \lambda_1 q_1 q_1^\top + \lambda_2 q_2 q_2^\top$.
  \item Compute $A^{10}$ using the spectral decomposition.
  \item Determine whether $A$ is positive definite.
\end{enumerate}
\end{exercise}

\begin{exercise}[PCA Application]
\label{exer:pca}
\textnormal{(\(**\))} \\
Consider centered data points $x_1 = (2, 1)^\top$, $x_2 = (-1, 1)^\top$,
$x_3 = (-1, -2)^\top$ in $\R^2$.
\begin{enumerate}[nosep,label=(\alph*)]
  \item Compute the sample covariance matrix
        $S = \frac{1}{3}\sum_{i=1}^3 x_i x_i^\top$.
  \item Find the eigenvalues and eigenvectors of $S$.
  \item Identify the first principal component direction and compute the
        projection of each data point onto it.
  \item What fraction of the total variance is captured by the first
        principal component?
\end{enumerate}
\end{exercise}

\begin{exercise}[Courant--Fischer Application]
\label{exer:courant_fischer}
\textnormal{(\(**\))} \\
Let $A \in \R^{n \times n}$ be symmetric with eigenvalues
$\lambda_1 \leq \cdots \leq \lambda_n$.
\begin{enumerate}[nosep,label=(\alph*)]
  \item Use the Courant--Fischer theorem to prove the \emph{interlacing
        inequality}: if $B$ is the $(n-1) \times (n-1)$ principal submatrix
        obtained by deleting the last row and column of $A$, and $B$ has
        eigenvalues $\mu_1 \leq \cdots \leq \mu_{n-1}$, then
        $\lambda_1 \leq \mu_1 \leq \lambda_2 \leq \mu_2 \leq \cdots
        \leq \mu_{n-1} \leq \lambda_n$.
  \item Apply this result: if $A \in \R^{3 \times 3}$ is symmetric with
        eigenvalues $1, 3, 5$, what are the possible eigenvalues of the
        $2 \times 2$ upper-left submatrix?
\end{enumerate}
\textit{Hint:} For (a), embed $(n-1)$-dimensional subspaces of
$\R^{n-1}$ into $\R^n$ via $x \mapsto (x, 0)$.
\end{exercise}

\begin{exercise}[Positive Definiteness and Convexity]
\label{exer:pd_proof}
\textnormal{(\(***\))} \\
\begin{enumerate}[nosep,label=(\alph*)]
  \item Let $A, B \in \R^{n \times n}$ be symmetric with $A \succeq B$.
        Use the Courant--Fischer theorem to prove $\lambda_k(A) \geq \lambda_k(B)$
        for each $k$.
  \item Let $A_0, A_1 \succ 0$ and $A(t) = (1-t)A_0 + tA_1$ for
        $t \in [0,1]$.  Show that $A(t) \succ 0$ for all $t$ (the positive
        definite cone is convex).
\end{enumerate}
\end{exercise}

%% ============================================================
%% SECTION 9: SUMMARY
%% ============================================================
\section{Summary}
\label{sec:summary}

\begin{itemize}
  \item An \emph{inner product} on a real vector space is a symmetric,
        bilinear, positive definite function $\langle \cdot, \cdot \rangle$.
        It induces a norm $\|v\| = \sqrt{\langle v, v \rangle}$ and a
        distance $d(u,v) = \|u - v\|$.

  \item The \emph{Cauchy--Schwarz inequality}
        $|\langle u, v \rangle| \leq \|u\| \cdot \|v\|$ is the most
        fundamental inequality in inner product spaces, with equality iff
        $u$ and $v$ are linearly dependent.

  \item The \emph{Gram--Schmidt process} converts any linearly independent
        list into an orthonormal list with the same successive spans.
        Every finite-dimensional inner product space has an orthonormal
        basis.

  \item For any subspace $U$, the space decomposes as
        $V = U \oplus U^\perp$.  The \emph{orthogonal projection}
        $\proj_U(v) = \sum_j \langle v, e_j \rangle e_j$ gives the unique
        closest point in $U$ to $v$ (best approximation theorem).

  \item \emph{Bessel's inequality}:
        $\sum_j |\langle v, e_j \rangle|^2 \leq \|v\|^2$ for any
        orthonormal list, with equality iff $v$ lies in the span.

  \item The \emph{spectral theorem} states that every real symmetric
        matrix has an orthonormal basis of eigenvectors with real
        eigenvalues: $A = Q\Lambda Q^\top$ with $Q$ orthogonal.

  \item A symmetric matrix is \emph{positive definite} iff all eigenvalues
        are positive, iff $x^\top Ax > 0$ for all $x \neq 0$, iff it admits
        a Cholesky factorization $A = LL^\top$.

  \item The \emph{Rayleigh quotient} $R_A(x) = x^\top Ax / x^\top x$
        satisfies $\lambda_{\min} \leq R_A(x) \leq \lambda_{\max}$.  The
        \emph{Courant--Fischer theorem} gives a minimax characterization
        of every eigenvalue, enabling eigenvalue comparison between matrices
        and submatrices.
\end{itemize}

%% ============================================================
%% SECTION 10: FURTHER READING
%% ============================================================
\section{Further Reading}
\label{sec:further_reading}

\begin{itemize}
  \item \textbf{Lesson 8} (SVD and Applications) extends the spectral
        theorem to non-symmetric and non-square matrices via the singular
        value decomposition $A = U\Sigma V^\top$.  The SVD is the most
        widely used matrix factorization in machine learning.

  \item \textbf{Phase 01, Functional Analysis} (Lessons on Hilbert spaces)
        generalizes the finite-dimensional inner product theory to
        infinite-dimensional spaces, with applications to kernel methods
        and reproducing kernel Hilbert spaces (RKHS).

  \item The complex analogue of the spectral theorem (for Hermitian
        matrices) is developed in \citet[Chapter~4]{horn2013}, along with
        Weyl's eigenvalue perturbation inequalities.

  \item \citet[Chapter~6]{axler2015} presents inner product spaces
        and the spectral theorem in the coordinate-free operator-theoretic
        framework, avoiding matrices until the final step.

  \item For the Cholesky factorization and its role in numerical linear
        algebra, see \citet[Chapter~6]{strang2016}.  The factorization
        is the computational backbone of Gaussian process inference in ML,
        where the kernel matrix is symmetric positive definite.

  \item \citet{petersen2012} provides a comprehensive collection of matrix
        identities, including derivative formulas for determinants, traces,
        and matrix inverses, which are essential for implementing gradient
        computations in neural networks.
\end{itemize}

\bibliography{linear-algebra-refs}
\end{document}
